\section{Background}

\label{sec:background}
% ══════════════════════════════════════════════════════════════════════════════

\subsection{Conservation Finance Landscape}

Global conservation finance flows through four main channels
\citep{deutz2020financing}:

\begin{table}[H]
\centering
\caption{Conservation finance by source (2019 estimates).}
\label{tab:finance}
\begin{tabular}{lcc}
\toprule
\textbf{Source} & \textbf{Annual flow (B\$)} & \textbf{Share} \\
\midrule
Government budgets     & 78--82  & 56\% \\
Biodiversity offsets   & 6--9    & 6\% \\
Philanthropic          & 4--10   & 5\% \\
Private/commercial     & 14--26  & 14\% \\
Natural infrastructure & 24--27  & 19\% \\
\midrule
\textbf{Total}         & 126--154 & 100\% \\
\bottomrule
\end{tabular}
\end{table}

The gap between current flows and estimated needs (\$700--1{,}000B/year
\citep{paulson2020financing}) implies that private capital mobilization
must increase by an order of magnitude.

\subsection{Impact Bonds: From Social to Environmental}

Social Impact Bonds (SIBs) emerged in 2010 with the Peterborough Prison
SIB \citep{liebman2011social}. Environmental adaptations include the DC
Water EIB (2016, \$25M for green infrastructure
\citep{nicola2013quantifying}), the Rhino Impact Bond (2020, \$45M for
black rhino conservation \citep{uct2020rhino}), and the Forest
Resilience Bond (2018, \$4M for wildfire risk reduction
\citep{blue2019forest}). These instruments demonstrate feasibility but
remain bespoke, with structuring costs of \$1--5M per issuance.

\subsection{Blockchain in Conservation}

Blockchain applications in conservation include supply chain
traceability \citep{howson2019crypto}, carbon credit registries
\citep{schletz2020blockchain}, and wildlife trade monitoring
\citep{dalton2020blockchain}. The integration of DeFi primitives with
conservation outcomes---the approach taken by BIBs---is novel.

\subsection{Biodiversity Credit Markets}

The nascent biodiversity credit market draws on lessons from carbon
markets. The TNFD framework \citep{tnfd2023framework} provides reporting
standards. BIBs differ from biodiversity credits by embedding outcomes in
financial instruments rather than creating offset commodities.


% ══════════════════════════════════════════════════════════════════════════════
